% 80_comparison：第 8 部，與 W06、W09、W10、W12 指定文獻之對照，以及釋字第 499 號。
% 頁碼核對：Michelman 415、421（PDF 第 19、25 頁）與 Post 437（PDF 第 10 頁）、441（PDF 第 14 頁，含註 45 之表）已對頁圖。
% Waldron 1368 頁與註 56 自 readings/W06 文字層讀過；W06 導讀問答 Teile/40_assumptions.tex 所列頁碼相同。
% Kramer 621--623、630、634 頁依 W09 預習 Teile/30_Kramer.tex 所載，並對 readings/W09 文字層。
% Goldstein 僅讀 391--392 頁（readings/W12 文字層），其論證本篇不寫。
% 釋字第 499 號：理由書自 phenom-court-data constitutionalCourt-reasoning-fulltext.json 取，依空行分為 12 段，段次與司法院大法官網站所示相同；
%   解釋文（五項）自 cons.judicial.gov.tw 該號解釋頁（id=310680）取得，該 JSON 不收解釋文。
% 第五次修憲經過：李炳南計畫報告為無文字層之掃描本，已辨讀第二章（library/text/李炳南_第五六次修憲之研究_第二章.txt，PDF 頁＝報告頁＋3），引用頁為〈計畫前言〉第 2 頁與第二章第 39--40、61 頁；修憲日期與內容以釋字第 499 號理由書為準（報告第 45 頁三讀日期誤作 9 月 3 日，不採）。
% Rawls《Political Liberalism》、Dworkin 等本機無藏本，一律轉引。
\chapter{前後週對照與中華民國法}
\label{ch:80_comparison}

\begin{qaquestion}
Michelman\ix{persons}{michelman@Michelman, Frank I.} 與 Post\ix{persons}{post@Post, Robert} 兩篇與前後各週指定文獻，以及中華民國司法院大法官解釋\ix{terms}{08@\zhlidxcat{憲法解釋與司法自制}!解釋}，各有何可對照之處？
\begin{enumerate}
\tightlist
\item W06 Waldron\ix{persons}{waldron@Waldron, Jeremy} 的第四項假設，與 Michelman 所稱合理的詮釋多元\ix{terms}{02@\zhlidxcat{正當性與合理歧見}!合理的詮釋多元}、Post 所引 Rawls\ix{persons}{rawls@Rawls, John} 的合理歧見\ix{terms}{02@\zhlidxcat{正當性與合理歧見}!合理歧見}，各自指向什麼？
\item W09 Kramer\ix{persons}{kramer@Kramer, Larry D.} 的人民憲政主義\ix{terms}{04@\zhlidxcat{違憲審查權限之歸屬}!人民憲政主義}與 Post 的人民主權\ix{terms}{01@\zhlidxcat{民主理論}!人民主權}，所稱的人民各有何種作為？
\item W10 Fallon\ix{persons}{fallon@Fallon, Richard H., Jr.} 的三種正當性\ix{terms}{02@\zhlidxcat{正當性與合理歧見}!正當性}與 Law\ix{persons}{law@Law, David S.}，與 Post 所說民主\ix{terms}{01@\zhlidxcat{民主理論}!民主}自治\ix{terms}{01@\zhlidxcat{民主理論}!自治}以外的正當性來源，如何對照？
\item W12 Goldstein\ix{persons}{goldstein@Goldstein, Leslie F.} 一文的題目，與本週兩篇有何銜接？
\item 釋字第 499 號\ix{cases}{01@中華民國!01@司法院大法官!0499@釋字第 499 號}解釋宣告第五次修憲之增修條文失其效力，與 Post 所舉依法定程序通過而反民主的修憲例，如何對照？
\end{enumerate}
\end{qaquestion}

\section{W06 Waldron：合理歧見}
\label{sec:80_comparison-1}

Waldron\ix{persons}{waldron@Waldron, Jeremy} 在 1368 頁寫道，多數情形下關於權利的合理歧見\ix{terms}{02@\zhlidxcat{正當性與合理歧見}!合理歧見}是「pursued reasonably and in good faith」〔出於合理與善意〕，並稱所假設者近於 Rawls 的 burdens of judgment〔判斷之負擔\ix{terms}{02@\zhlidxcat{正當性與合理歧見}!判斷的負擔}〕，差別在適用範圍：Rawls 在 right〔正當〕的議題上遲疑，Waldron 將其同時適用於 right 與 good〔善〕。\footcite[1368 \& n.56]{waldron2006}承認這種歧見與道德相對主義無關，因為歧見的存在與權利及憲政主義\ix{terms}{03@\zhlidxcat{憲政主義與憲法之效力基礎}!憲政主義}原則另有真理一事相容。\footcite[1368]{waldron2006}

Michelman 在 421 頁稱自己所設為「the fact of reasonable interpretive pluralism」〔合理的詮釋多元\ix{terms}{02@\zhlidxcat{正當性與合理歧見}!合理的詮釋多元}的事實〕，內容是自由國家居民之間，對正義所要求之根本法\ix{terms}{03@\zhlidxcat{憲政主義與憲法之效力基礎}!根本法}解釋，存在「inevitable irresolvable uncertainty and irreparable reasonable disagreement」〔無可避免而無從解決的不確定，與無從化解之合理歧見\ix{terms}{02@\zhlidxcat{正當性與合理歧見}!無從化解之合理歧見}〕。他表明並非嚴格依循 Rawls：Rawls 所設為整全性學說\ix{terms}{02@\zhlidxcat{正當性與合理歧見}!整全性學說}（comprehensive views）的合理多元\ix{terms}{02@\zhlidxcat{正當性與合理歧見}!合理多元}，Michelman 所設為參與者對自由憲政原則之解釋的合理多元，且認為合理多元連同判斷之負擔「cuts deeper」〔切入更深〕。\footcite[421 \& nn.65--67]{michelman1998}422 頁接著說，合理多元並未使此一領域的真理在哲學上不可得，使之不可得的是政治上：參與者各自相信有真理，卻無從就其內容達成一致，也無從把所有相反立場都斥為不合理。\footcite[422]{michelman1998}Waldron 在 1368 頁同樣把歧見之承認與真理之有無分開。

兩人對這項前提的使用不同。Waldron 以之為第四項假設，W06 的問題是在這項假設之下，由法院還是立法機關作最後決定；\footcite[1366--1369]{waldron2006}Michelman 則用它說明解釋問題終須有人決定，「Whoever does not participate in the decision, for whatever reason, is to that extent governed by those who do」〔不論基於何種理由而未參與決定者，在此範圍內受參與決定者所治理〕。\footcite[415]{michelman1998}Post 引用同一組 Rawls 用語，作用又另有所在：433 頁以整全性\ix{terms}{08@\zhlidxcat{憲法解釋與司法自制}!整全性}學說的多樣性說明，大部分人口不接受任何大法官的判斷為正當；436 頁以「irreparable reasonable disagreement」說明，民主正當性\ix{terms}{02@\zhlidxcat{正當性與合理歧見}!民主正當性}獨立於立法結果，因此與歧見相容，Michelman 模型則使歧見在一切情形下都否定民主正當性。\footcite[433 n.19, 436 \& n.28]{post1998}

\begin{longtable}{@{}L{0.14\linewidth}L{0.28\linewidth}L{0.27\linewidth}L{0.20\linewidth}@{}}
\caption{三篇文獻關於合理歧見之用語。}\\
\toprule
文獻 & 用語 & 所指對象 & 頁次 \\
\midrule
\endfirsthead
\toprule
文獻 & 用語 & 所指對象 & 頁次 \\
\midrule
\endhead
Waldron & burdens of judgment〔判斷之負擔〕，適用於 right 與 good & 權利之內容與核心適用 & 1368 及註 56 \\
Michelman & reasonable interpretive pluralism〔合理的詮釋多元〕 & 自由憲政原則之解釋 & 421--422 \\
Post & irreparable reasonable disagreement〔無從化解之合理歧見〕，引 Rawls & 現代異質國家之社會歧見 & 436 及註 28 \\
Post & comprehensive doctrines〔整全性學說〕，引 Rawls & 人口對大法官判斷之接受 & 433 及註 19 \\
\bottomrule
\end{longtable}

\section{W09 Kramer：人民憲政主義與人民主權}
\label{sec:80_comparison-2}

Kramer\ix{persons}{kramer@Kramer, Larry D.} 621 頁的摘要把共和派（Jeffersonian）的立場寫成：詮釋憲法的首要權威在人民，不在法院，即 popular constitutionalism〔人民憲政主義\ix{terms}{04@\zhlidxcat{違憲審查權限之歸屬}!人民憲政主義}〕。622 至 623 頁進一步說，在這套理論中憲法之於政府各部門，如同法律之於個人；官員對憲法的詮釋與人民對普通法的詮釋同樣不具權威，受人民自己更高權威的直接矯正。\footcite[621--623]{kramer2012}依此，人民除制定憲法之外，亦實際詮釋憲法。

Post 對 popular sovereignty〔人民主權〕的界定較窄。437 頁將其界定為國家服從人民意志\ix{terms}{01@\zhlidxcat{民主理論}!人民意志}，而人民意志依多數決或憲法第五條修憲程序這類程序標準認定；\footcite[437]{post1998}440 頁稱這些標準發揮承認規則\ix{terms}{03@\zhlidxcat{憲政主義與憲法之效力基礎}!承認規則}（rule of recognition）的作用，把公眾意見\ix{terms}{01@\zhlidxcat{民主理論}!公眾意見}轉為實證法，憲法解釋因此始終可由人民依該程序之意志修正，而司法對國族精神\ix{terms}{03@\zhlidxcat{憲政主義與憲法之效力基礎}!國族精神}的闡述，在人民主權的縫隙之間生長。\footcite[440]{post1998}Post 的人民意志經由修憲程序表達，Kramer 的人民則握有詮釋權威。

三篇對詮釋之份額的安排如下。

\begin{longtable}{@{}L{0.16\linewidth}L{0.46\linewidth}L{0.28\linewidth}@{}}
\caption{三篇文獻對憲法詮釋權之歸屬。}\\
\toprule
文獻 & 內容 & 頁次 \\
\midrule
\endfirsthead
\toprule
文獻 & 內容 & 頁次 \\
\midrule
\endhead
Michelman & 自治論之下，解決基本而有爭議的政治道德問題者，應是人民依民主方式行事；憲法解釋之主要工作不能置於民主範圍之外 & 415 \\
Post & 憲法解釋始終受修憲程序所定之人民意志修正；司法闡述國族精神，則在人民主權的縫隙之間 & 440 \\
Kramer & 共和派理論以人民為首要詮釋者，官員之詮釋受人民直接矯正 & 621--623（W09） \\
\bottomrule
\end{longtable}

Kramer 與 Post 都寫到 1930 年代新政時期。Kramer 630 頁引 1937 年憲法日演說中 Roosevelt 之語，稱憲法是「a layman's document, not a lawyer's contract」〔外行人的文件，不是律師的契約〕，其意義由人民與國會決定，Roosevelt 此後藉任命權鞏固勝利。\footcite[630 \& n.33]{kramer2012}Post 441 至 442 頁則把新政危機作為法院將世代（或階級）理想強加於敵意人口的例子，並列舉法院堅持特定國族精神時可能遭受的政治報復：公民不服從、敵意的總統任命、修憲。\footcite[441--442]{post1998}兩人所述的修憲與任命，皆屬人民意志對法院的回應；Kramer 把它放在詮釋權限之歸屬，Post 放在法院闡述國族精神的風險。

\section{W10 Fallon 與 Law：正當性}
\label{sec:80_comparison-3}

Fallon\ix{persons}{fallon@Fallon, Richard H., Jr.} 21 至 22 頁區分三種正當性：社會學意義的正當性取決於公眾對政府、機關或裁判的實際態度，道德意義的正當性涉及人們不論他人如何想都應當持有的態度，法律意義的正當性是法律概念。\footcite[21--22]{fallon2018}Law 779 頁批評把正當性當成可累積與耗盡的蓄水池，793 頁稱此說迴避了服從的動機；800 至 801 頁所列新問題，如人民為何服從法院、政府為何遵從裁判，屬服從問題。\footcite[779, 793, 800--801]{law2009}

Post 使用 legitimacy 的三處如下。435 頁，Post 承認 Michelman 的兩項結構條件大致可建立對抗制\ix{terms}{07@\zhlidxcat{法院之組織與功能}!對抗制}（adversary system）裁判之正當性，Fuller\ix{persons}{fuller@Fuller, Lon L.} 等人亦如此相信，但無人主張它足以使對抗制成為民主自治的結構。436 頁，這兩項條件至多只能建立 political legitimacy〔政治正當性\ix{terms}{02@\zhlidxcat{正當性與合理歧見}!政治正當性}〕，無從產生民主自治的價值。439 頁，憲法對民主價值的解釋不因欠缺民主自治而失其正當性，因為正當性另有許多民主自決以外的來源：民主之實踐嵌在非民主的社會組織形式之中，社群\ix{terms}{01@\zhlidxcat{民主理論}!社群}（community）規範藉社會化與文化再生產而延續。\footcite[435--436, 439]{post1998}

Post 依正當性之來源分為民主自決與社群規範\ix{terms}{01@\zhlidxcat{民主理論}!社群規範}等；Fallon 依所指涉的概念分為事實態度、應然判斷與法律概念。Post 在這幾頁未把其所說的 political legitimacy 與 legitimacy 歸入 Fallon 的任何一類。可對照之處有二：435 頁所述 1920 年代對 Taft 法院的攻擊，描述公眾對法院的態度；\footcite[435]{post1998}440 頁的承認規則\ix{terms}{03@\zhlidxcat{憲政主義與憲法之效力基礎}!承認規則}把人民意志\ix{terms}{01@\zhlidxcat{民主理論}!人民意志}轉為實證法，涉及法律上之效力。

\section{W12 Goldstein：guardian democracy 與 representative democracy}
\label{sec:80_comparison-4}

Goldstein\ix{persons}{goldstein@Goldstein, Leslie F.} 的題目為〈Judicial Review and Democratic Theory: Guardian Democracy vs. Representative Democracy〉，391 至 392 頁自述其問題是：在「extra-textualist」等學派與 Marbury v. Madison 以來的傳統之間，兩種違憲審查\ix{terms}{04@\zhlidxcat{違憲審查權限之歸屬}!違憲審查}理論各為人民或其民選代表留下何種位置，即是否仍有「government by consent of the governed」〔經被治者同意之政府〕的空間。\footcite[391--392]{goldstein1987}本篇不預述其論證。

與本週銜接之處限於用語。Post 437 頁引 Brown v. Hartlage\ix{cases}{02@美國!01@聯邦最高法院!1982Brown v. Hartlage（1982）@Brown v. Hartlage（1982）} 的措辭，稱第一修正案為「the guardian of our democracy」〔我們民主的守護者〕，其所指為第一修正案所保障之言論自由，並未指法官。\footcite[437 \& n.30]{post1998}Goldstein 391 頁以「a bevy of Platonic Guardians」〔一群柏拉圖式的守護者〕指稱 Learned Hand 所批評的統治。\footcite[391]{goldstein1987}Post 432 頁引 Michelman 424 至 426 頁，以公民出於尊重\ix{terms}{02@\zhlidxcat{正當性與合理歧見}!出於尊重}（respect）而遵守根本法為標準；\footcite[432 \& nn.12--14]{post1998}Goldstein 392 頁以被治者之同意為問題。下週可問這兩項標準是否相同。

\section{中華民國司法院大法官釋字第 499 號解釋}
\label{sec:80_comparison-5}

\subsection{Post 之修憲例}

Post 在 430 頁設想美國聯邦憲法修正案經國會兩院各三分之二提出、四分之三州議會批准，廢止增修條文第 26 條並明定未滿 21 歲者不得投票。431 頁假設十八歲公民據以訴請投票權，法官若判原告勝訴，即屬未能理解憲法裁判之實踐；法官縱認該修正案反民主，仍須依第五條程序所量度之人民意志\ix{terms}{01@\zhlidxcat{民主理論}!人民意志}執行，此時「popular sovereignty manifestly trumps constitutionalism」〔人民主權明顯勝過憲政主義〕。\footcite[430--431]{post1998}437 頁以此為人民主權而不必然是民主之例，另設更極端之例：人民若依法定程序制定修正案，廢除投票並授予聯邦官員終身及世襲任期，該修正案是人民主權的例子，卻明顯反民主。\footcite[437]{post1998}431 頁註 10 稱，明確授權總統終身任期之修正案，會顯現民主與人民主權之間的緊張。\footcite[431 n.10]{post1998}Post 441 頁註 45 的表以民主之為規範理想或人民主權，對憲政主義之為國族精神\ix{terms}{03@\zhlidxcat{憲政主義與憲法之效力基礎}!國族精神}或實證法，共分四格；實證法與民主理想的關係屬第 2 格，所舉為第五條修憲程序這類實證憲法要求。\footcite[441 n.45]{post1998}

\subsection{第五次修憲與解釋之作成}

釋字第 499 號\ix{cases}{01@中華民國!01@司法院大法官!0499@釋字第 499 號}解釋的聲請人為立法委員，聲請標的是 1999 年 9 月 15 日公布之憲法增修條文，其中國民大會 1999 年 9 月 4 日三讀通過之部分，二讀會與三讀會皆採無記名投票。\footcites[理由書第 1 段]{jy499}[39--40]{li2001}依李炳南的研究計畫報告，第五、六次修憲之成果主要來自國民大會制度內部所發動之提案，立法院史無前例地提案聲請大法官解釋部分增修條文無效。\footcite[計畫前言 2；61]{li2001}司法院大法官於 2000 年 3 月 24 日作成解釋。

\subsection{程序與實體之審查}

解釋文第一項與理由書第 5 段稱，修改憲法是最直接體現國民主權之行為，國民主權之行使「應公開透明以滿足理性溝通之條件，方能賦予憲政國家之正當性基礎」。\footcite[解釋文第一項；理由書第 5 段]{jy499}修憲如有重大明顯瑕疵，即不生其應有之效力；二讀與三讀採無記名投票，已達重大明顯之程度。\footcite[理由書第 6 段]{jy499}對國民大會所稱修憲程序不受司法審查，理由書第 7 段逐一回應，其中對美國部分指出，Coleman v. Miller\ix{cases}{02@美國!01@聯邦最高法院!1939Coleman v. Miller（1939）@Coleman v. Miller（1939）}, 307 U.S. 433 (1939) 一案並未形成修憲程序均屬政治問題之多數意見，並引 Uhler v. AFL-CIO, 468 U.S. 1310 (1984) 為證。\footcite[理由書第 7 段]{jy499}

實體部分，理由書第 9 段稱：「我國憲法雖未明定不可變更之條款」，但第 1 條之民主共和國原則\ix{statutes}{01@中華民國!01@中華民國憲法!001@第 1 條}、第 2 條之國民主權原則\ix{statutes}{01@中華民國!01@中華民國憲法!002@第 2 條}、人民權利保障與權力分立制衡，具有本質之重要性，聽任修改條文予以變更，則「憲法上整體規範秩序將形同破毀，此等修改之條文則失其應有之正當性」。涉及自由民主憲政秩序之違反者，「已悖離國民之付託」，應受憲法所設置其他權力部門之制約。\footcite[解釋文第二項；理由書第 9 段]{jy499}

增修條文第 1 條使國民大會代表依立法委員選舉各政黨得票比例分配席次，理由書第 10 段稱，其本身未曾辦理選舉，產生之代表「已不具民意代表身分，充其量為各政黨指派之代表」，與民主之憲政秩序有違；該段並回應國民大會以兩院制國家之上議院為例的主張，指出任命、世襲制度「已為現代大多數民主國家所不取」。\footcite[理由書第 10 段]{jy499}增修條文將第三屆國民大會代表任期延長二年又四十二天、第四屆立法委員任期延長五個月；理由書第 11 段稱，代議民主之正當性在於遵守與選民之約定，任期屆滿除有不能改選之正當理由外應即改選，所謂正當理由須與釋字第 31 號\ix{cases}{01@中華民國!01@司法院大法官!0031@釋字第 31 號}解釋所指「國家發生重大變故，事實上不能依法辦理次屆選舉」相當，無故延任者「難謂具有正當性」。\footcite[理由書第 11 段]{jy499}

解釋文第五項宣告第 1 條、第 4 條、第 9 條及第 10 條之修正自解釋公布之日起失其效力，1997 年 7 月 21 日修正公布之原增修條文繼續適用；第 9 條與第 10 條內容無可議，因通過程序違背修憲正當程序而一併失效。\footcite[解釋文第五項；理由書第 12 段]{jy499}

\subsection{與 Post 修憲例之對照}

\begin{longtable}{@{}L{0.15\linewidth}L{0.38\linewidth}L{0.39\linewidth}@{}}
\caption{Post 之假設修正案與第五次修憲增修條文。}\\
\toprule
項目 & Post 之假設修正案 & 釋字第 499 號\ix{cases}{01@中華民國!01@司法院大法官!0499@釋字第 499 號}所審查之增修條文 \\
\midrule
\endfirsthead
\toprule
項目 & Post 之假設修正案 & 釋字第 499 號所審查之增修條文 \\
\midrule
\endhead
通過程序 & 國會兩院各三分之二提出、四分之三州議會批准（430）；437 頁之例同為依法定程序制定 & 二讀、三讀採無記名投票，大法官認為瑕疵重大明顯（解釋文第一項；理由書第 6 段） \\
選舉 & 限 21 歲以上始得投票（430）；廢除投票（437） & 國民大會代表依立法委員選舉得票比例分配席次，本身未經選舉（理由書第 10 段） \\
任期 & 聯邦官員終身及世襲（437）；總統終身（431 註 10） & 第三屆國民大會代表延長二年又四十二天，第四屆立法委員延長五個月（理由書第 11 段） \\
法院之處理 & 法官依第五條所量度之人民意志\ix{terms}{01@\zhlidxcat{民主理論}!人民意志}執行修正案，人民主權勝過憲政主義（431） & 宣告第 1、4、9、10 條失其效力，原增修條文繼續適用（解釋文第五項） \\
所用概念 & popular sovereignty 與 democracy（437--438） & 國民主權原則（憲法第 2 條）、自由民主憲政秩序（解釋文第二項） \\
\bottomrule
\end{longtable}

Post 的例所取消的是投票，並授予終身與世襲任期，增修條文所涉者為代表產生方式與任期延長若干月，兩者程度不同。釋字第 499 號與 Post 的概念有兩處可對照。第一，Post 把人民主權界定為依程序標準認定的人民意志；\footcite[437]{post1998}大法官在理由書第 5 段，將國民主權之行使繫於公開透明與理性溝通的條件，又在第 11 段稱民意代表之權限應直接源自國民之授權。第二，Post 設想的法官在 Article V 程序下執行修正案；大法官則在解釋文第二項以本質重要之憲法原則為限制，審查依修憲程序制定之條文，並於理由書第 9 段稱此為憲法自我防衛之機制。\footcite[解釋文第二項；理由書第 5、9、11 段]{jy499}釋字第 499 號把民主共和國原則與國民主權原則並列為本質重要之憲法原則，Post 441 頁註 45 的表則把人民主權與民主之規範理想分列兩列。

\section{課前可提的問題}
\label{sec:80_comparison-6}

\begin{enumerate}
\tightlist
\item Waldron\ix{persons}{waldron@Waldron, Jeremy} 1368 頁假設權利之分歧出於合理與善意，Post 436 至 437 頁稱民主正當性\ix{terms}{02@\zhlidxcat{正當性與合理歧見}!民主正當性}獨立於立法結果，因此與無從化解之合理歧見\ix{terms}{02@\zhlidxcat{正當性與合理歧見}!無從化解之合理歧見}相容。Waldron 的假設與 Post 的這項獨立性，能否同時成立？
\item Post 437 與 440 頁以修憲程序界定人民主權，Kramer\ix{persons}{kramer@Kramer, Larry D.} 621 至 623 頁的共和派以人民為首要詮釋者。後者的人民所作的詮釋，在 Post 的定義下屬於人民主權還是民主？
\item Post 439 頁稱正當性另有民主自決以外的來源，Fallon\ix{persons}{fallon@Fallon, Richard H., Jr.} 21 至 22 頁將正當性分為社會學、道德與法律三種意義。Post 所舉社群規範\ix{terms}{01@\zhlidxcat{民主理論}!社群規範}之延續，屬於其中哪一種？
\item 釋字第 499 號\ix{cases}{01@中華民國!01@司法院大法官!0499@釋字第 499 號}理由書第 5、11 段把國民主權繫於公開透明與選民授權，Post 437 頁把人民主權界定為依程序標準認定的人民意志\ix{terms}{01@\zhlidxcat{民主理論}!人民意志}。依 Post 的定義，大法官所用的國民主權，是否已含有 Post 歸給民主的實質內容？
\item Post 437 頁引「the guardian of our democracy」指第一修正案，Goldstein\ix{persons}{goldstein@Goldstein, Leslie F.} 391 頁以 Platonic Guardians 指稱法官之治。下週 Goldstein 所稱的 guardian democracy，是 Post 所指的哪一種民主？
\end{enumerate}
